
\documentclass{article}
\usepackage{colortbl}
\usepackage{makecell}
\usepackage{multirow}
\usepackage{supertabular}

\begin{document}

\newcounter{utterance}

\centering \large Interaction Transcript for game `function\_detective', experiment `two\_numbers\_passive\_labeled\_pairs', episode 291 with Qwen3.6{-}35B{-}A3B{-}FP8{-}without{-}reasoning.
\vspace{24pt}

{ \footnotesize  \setcounter{utterance}{1}
\setlength{\tabcolsep}{0pt}
\begin{supertabular}{c@{$\;$}|p{.15\linewidth}@{}p{.15\linewidth}p{.15\linewidth}p{.15\linewidth}p{.15\linewidth}p{.15\linewidth}}
    \# & \multicolumn{2}{c}{Player} && \multicolumn{2}{c}{Game Master} \\
    \hline

    \theutterance \stepcounter{utterance}  
    & & & \multicolumn{4}{p{0.6\linewidth}}{
        \cellcolor[rgb]{0.9,0.9,0.9}{
            \makecell[{{p{\linewidth}}}]{
                \texttt{\tiny{[P1$\langle$GM]}}
                \texttt{I have a hidden Python function for you to figure out.} \\
\\ 
\texttt{Your job is to infer what it does and then submit a correct Python implementation.} \\
\\ 
\texttt{Function signature:} \\
\texttt{(a: int, b: int) {-}> int} \\
\\ 
\texttt{You have 15 turns total.} \\
\\ 
\texttt{I reveal candidate input{-}output pairs one at a time and tell you whether each one is correct.} \\
\\ 
\texttt{You must make your best SOLVE attempt no later than the second{-}to{-}last turn.} \\
\texttt{Do not wait until the very end.} \\
\\ 
\texttt{You can use these actions:} \\
\texttt{{-} JSON format described at the end of the prompt} \\
\texttt{{-} SOLVE: ```python ... ```} \\
\\ 
\texttt{Do not add import statements in your guess.} \\
\texttt{Anything you might need is already available, so just write the function body in hypothesis and return the answer.} \\
\\ 
\texttt{Preloaded examples:} \\
\texttt{1. ((({-}99, 10), 723), True)} \\
\\ 
\texttt{Scoring:} \\
\texttt{{-} correctness is required for any non{-}zero score} \\
\texttt{{-} in interactive modes, solving with fewer information requests is better} \\
\\ 
\texttt{Input formatting reminders:} \\
\texttt{{-} strings must be quoted} \\
\texttt{{-} booleans must be True/False} \\
\texttt{{-} use commas between multiple values} \\
\texttt{{-} do not write any extra action tags in the same message} \\
\\ 
\texttt{Format example:} \\
\texttt{You write:} \\
\texttt{\{'current\_mode': <probing/confirming/don't know; choose whichever best describes your current state in the task>, 'input\_rationale': <brief description of why you chose this hypothesis>, 'hypothesis': <either the exact string "unknown" if you have no guess yet OR your current best guess as runnable Python matching the given signature, with a one{-}line docstring stating what it does. Never write placeholder code, if you cannot commit to a guess, use "unknown". Match the given signature.} \\
\texttt{e.g.} \\
\texttt{def f(x):} \\
\texttt{	"""Returns the absolute value of x."""} \\
\texttt{	return x if x >= 0 else {-}x} \\
\\ 
\texttt{Use "unknown" if you have no guess.>, 'confidence\_score': <score between 1{-}100. 100 means highest confidence in your hypothesis>\}} \\
\\ 
\texttt{Output only one valid json per turn.} \\
            }
        }
    }
     \\ \\

    \theutterance \stepcounter{utterance}  
    & \multicolumn{4}{p{0.6\linewidth}}{
        \cellcolor[rgb]{0.561,0.737,0.561}{
            \makecell[{{p{\linewidth}}}]{
                \texttt{\tiny{[P1$\rangle$GM]}}
                \texttt{\{'current\_mode': 'probing', 'input\_rationale': 'I have one preloaded example: f({-}99, 10) = 723. I need to probe more to understand the function. Let me try some simple cases to see if it is a basic arithmetic operation or something more complex.', 'hypothesis': 'unknown', 'confidence\_score': 10\}} \\
            }
        }
    }
    & & \\ \\

    \theutterance \stepcounter{utterance}  
    & & & \multicolumn{4}{p{0.6\linewidth}}{
        \cellcolor[rgb]{0.9,0.9,0.9}{
            \makecell[{{p{\linewidth}}}]{
                \texttt{\tiny{[P1$\langle$GM]}}
                \texttt{Game Violation: ParseError: Response is not valid JSON.} \\
\texttt{Output ONLY:} \\
\texttt{One JSON with confidence\_score, current\_mode, input\_rationale, and hypothesis keys ...} \\
\texttt{OR} \\
\texttt{SOLVE: ```python ... ```} \\
            }
        }
    }
     \\ \\

    \theutterance \stepcounter{utterance}  
    & \multicolumn{4}{p{0.6\linewidth}}{
        \cellcolor[rgb]{0.561,0.737,0.561}{
            \makecell[{{p{\linewidth}}}]{
                \texttt{\tiny{[P1$\rangle$GM]}}
                \texttt{\{} \\
\texttt{  "current\_mode": "probing",} \\
\texttt{  "input\_rationale": "I have one preloaded example: f({-}99, 10) = 723. I need to probe more to understand the function. Let me try some simple cases to see if it is a basic arithmetic operation or something more complex.",} \\
\texttt{  "hypothesis": "unknown",} \\
\texttt{  "confidence\_score": 10} \\
\texttt{\}} \\
            }
        }
    }
    & & \\ \\

    \theutterance \stepcounter{utterance}  
    & & & \multicolumn{4}{p{0.6\linewidth}}{
        \cellcolor[rgb]{0.9,0.9,0.9}{
            \makecell[{{p{\linewidth}}}]{
                \texttt{\tiny{[P1$\langle$GM]}}
                \texttt{OUTPUT: ((({-}1, {-}99), {-}398), False)} \\
            }
        }
    }
     \\ \\

    \theutterance \stepcounter{utterance}  
    & \multicolumn{4}{p{0.6\linewidth}}{
        \cellcolor[rgb]{0.561,0.737,0.561}{
            \makecell[{{p{\linewidth}}}]{
                \texttt{\tiny{[P1$\rangle$GM]}}
                \texttt{\{} \\
\texttt{  "current\_mode": "probing",} \\
\texttt{  "input\_rationale": "I have two examples now: f({-}99, 10) = 723 and f({-}1, {-}99) = {-}398. Let's analyze these. \textbackslash{}n\textbackslash{}nExample 1: a={-}99, b=10, result=723.\textbackslash{}nExample 2: a={-}1, b={-}99, result={-}398.\textbackslash{}n\textbackslash{}nLet's look for patterns. \textbackslash{}n723 = {-}99 * 10 + 1722? No.\textbackslash{}n723 = {-}99 * ({-}10) + 1723? No.\textbackslash{}n\textbackslash{}nLet's try linear combination: c1*a + c2*b = result.\textbackslash{}n{-}99*c1 + 10*c2 = 723\textbackslash{}n{-}1*c1 {-} 99*c2 = {-}398\textbackslash{}n\textbackslash{}nFrom second eq: c1 = 398 {-} 99*c2.\textbackslash{}nSubstitute into first: {-}99*(398 {-} 99*c2) + 10*c2 = 723\textbackslash{}n{-}39402 + 9801*c2 + 10*c2 = 723\textbackslash{}n9811*c2 = 40125\textbackslash{}nc2 =} \\
            }
        }
    }
    & & \\ \\

    \theutterance \stepcounter{utterance}  
    & & & \multicolumn{4}{p{0.6\linewidth}}{
        \cellcolor[rgb]{0.9,0.9,0.9}{
            \makecell[{{p{\linewidth}}}]{
                \texttt{\tiny{[P1$\langle$GM]}}
                \texttt{Game Violation: ParseError: Response is not valid JSON.} \\
\texttt{Output ONLY:} \\
\texttt{One JSON with confidence\_score, current\_mode, input\_rationale, and hypothesis keys ...} \\
\texttt{OR} \\
\texttt{SOLVE: ```python ... ```} \\
            }
        }
    }
     \\ \\

    \theutterance \stepcounter{utterance}  
    & \multicolumn{4}{p{0.6\linewidth}}{
        \cellcolor[rgb]{0.561,0.737,0.561}{
            \makecell[{{p{\linewidth}}}]{
                \texttt{\tiny{[P1$\rangle$GM]}}
                \texttt{\{} \\
\texttt{  "current\_mode": "probing",} \\
\texttt{  "input\_rationale": "I have two examples: f({-}99, 10) = 723 and f({-}1, {-}99) = {-}398. The previous linear combination attempt didn't yield integer coefficients easily. Let's try more simple inputs to see if it's a standard operation like multiplication, addition, or something involving bitwise operations or string conversions. Let's try f(0, 0), f(1, 1), f(1, 0), f(0, 1) to get a baseline.",} \\
\texttt{  "hypothesis": "unknown",} \\
\texttt{  "confidence\_score": 20} \\
\texttt{\}} \\
            }
        }
    }
    & & \\ \\

    \theutterance \stepcounter{utterance}  
    & & & \multicolumn{4}{p{0.6\linewidth}}{
        \cellcolor[rgb]{0.9,0.9,0.9}{
            \makecell[{{p{\linewidth}}}]{
                \texttt{\tiny{[P1$\langle$GM]}}
                \texttt{OUTPUT: ((({-}99, {-}2), {-}501), True)} \\
            }
        }
    }
     \\ \\

    \theutterance \stepcounter{utterance}  
    & \multicolumn{4}{p{0.6\linewidth}}{
        \cellcolor[rgb]{0.561,0.737,0.561}{
            \makecell[{{p{\linewidth}}}]{
                \texttt{\tiny{[P1$\rangle$GM]}}
                \texttt{\{} \\
\texttt{  "current\_mode": "probing",} \\
\texttt{  "input\_rationale": "I have three examples now:\textbackslash{}n1. f({-}99, 10) = 723\textbackslash{}n2. f({-}1, {-}99) = {-}398\textbackslash{}n3. f({-}99, {-}2) = {-}501\textbackslash{}n\textbackslash{}nLet's look at the relationship between inputs and output.\textbackslash{}nExample 1: a={-}99, b=10, res=723.\textbackslash{}nExample 3: a={-}99, b={-}2, res={-}501.\textbackslash{}n\textbackslash{}nNotice that 'a' is the same in both ({-}99). Let's see how the result changes with 'b'.\textbackslash{}nWhen b goes from 10 to {-}2 (change of {-}12), the result goes from 723 to {-}501 (change of {-}1224).\textbackslash{}nSlope = {-}1224 / {-}12 = 102.\textbackslash{}n\textbackslash{}nSo, it seems like f(a, b) might be linear in b with a coefficient of 102.\textbackslash{}nLet's check: f(a, b) = 102 * b + C(a).\textbackslash{}nFor a={-}99, b=10: 102*10 + C({-}99) = 723 => 1020 + C} \\
            }
        }
    }
    & & \\ \\

    \theutterance \stepcounter{utterance}  
    & & & \multicolumn{4}{p{0.6\linewidth}}{
        \cellcolor[rgb]{0.9,0.9,0.9}{
            \makecell[{{p{\linewidth}}}]{
                \texttt{\tiny{[P1$\langle$GM]}}
                \texttt{Game Violation: ParseError: Response is not valid JSON.} \\
\texttt{Output ONLY:} \\
\texttt{One JSON with confidence\_score, current\_mode, input\_rationale, and hypothesis keys ...} \\
\texttt{OR} \\
\texttt{SOLVE: ```python ... ```} \\
            }
        }
    }
     \\ \\

    \theutterance \stepcounter{utterance}  
    & \multicolumn{4}{p{0.6\linewidth}}{
        \cellcolor[rgb]{0.561,0.737,0.561}{
            \makecell[{{p{\linewidth}}}]{
                \texttt{\tiny{[P1$\rangle$GM]}}
                \texttt{\{} \\
\texttt{  "current\_mode": "probing",} \\
\texttt{  "input\_rationale": "I have three examples: f({-}99, 10) = 723, f({-}1, {-}99) = {-}398, f({-}99, {-}2) = {-}501. The change in output with respect to b when a is constant ({-}99) suggests a coefficient of 102 for b. Specifically, 723 {-} ({-}501) = 1224, and 10 {-} ({-}2) = 12, so 1224/12 = 102. This suggests f(a, b) = 102*b + g(a). Let's test this hypothesis with the second example. If f({-}1, {-}99) = {-}398, then 102*({-}99) + g({-}1) = {-}398 => {-}10098 + g({-}1) = {-}398 => g({-}1) = 9700. Now let's check if this holds for a={-}99. f({-}99, 10) = 102*10 + g({-}99) = 723 => 1020 + g({-}99) = 723 => g({-}99) = {-}297} \\
            }
        }
    }
    & & \\ \\

    \theutterance \stepcounter{utterance}  
    & & & \multicolumn{4}{p{0.6\linewidth}}{
        \cellcolor[rgb]{0.9,0.9,0.9}{
            \makecell[{{p{\linewidth}}}]{
                \texttt{\tiny{[P1$\langle$GM]}}
                \texttt{Game Violation: ParseError: Response is not valid JSON.} \\
\texttt{Output ONLY:} \\
\texttt{One JSON with confidence\_score, current\_mode, input\_rationale, and hypothesis keys ...} \\
\texttt{OR} \\
\texttt{SOLVE: ```python ... ```} \\
            }
        }
    }
     \\ \\

    \theutterance \stepcounter{utterance}  
    & \multicolumn{4}{p{0.6\linewidth}}{
        \cellcolor[rgb]{0.561,0.737,0.561}{
            \makecell[{{p{\linewidth}}}]{
                \texttt{\tiny{[P1$\rangle$GM]}}
                \texttt{\{} \\
\texttt{  "current\_mode": "probing",} \\
\texttt{  "input\_rationale": "I have three examples: f({-}99, 10) = 723, f({-}1, {-}99) = {-}398, f({-}99, {-}2) = {-}501. The change in output with respect to b when a is constant ({-}99) suggests a coefficient of 102 for b. Specifically, 723 {-} ({-}501) = 1224, and 10 {-} ({-}2) = 12, so 1224/12 = 102. This suggests f(a, b) = 102*b + g(a). Let's test this hypothesis with the second example. If f({-}1, {-}99) = {-}398, then 102*({-}99) + g({-}1) = {-}398 => {-}10098 + g({-}1) = {-}398 => g({-}1) = 9700. Now let's check if this holds for a={-}99. f({-}99, 10) = 102*10 + g({-}99) = 723 => 1020 + g({-}99) = 723 => g({-}99) = {-}297} \\
            }
        }
    }
    & & \\ \\

    \theutterance \stepcounter{utterance}  
    & & & \multicolumn{4}{p{0.6\linewidth}}{
        \cellcolor[rgb]{0.9,0.9,0.9}{
            \makecell[{{p{\linewidth}}}]{
                \texttt{\tiny{[P1$\langle$GM]}}
                \texttt{Game Violation: ParseError: Response is not valid JSON.} \\
\texttt{Output ONLY:} \\
\texttt{One JSON with confidence\_score, current\_mode, input\_rationale, and hypothesis keys ...} \\
\texttt{OR} \\
\texttt{SOLVE: ```python ... ```} \\
            }
        }
    }
     \\ \\

    \theutterance \stepcounter{utterance}  
    & \multicolumn{4}{p{0.6\linewidth}}{
        \cellcolor[rgb]{0.561,0.737,0.561}{
            \makecell[{{p{\linewidth}}}]{
                \texttt{\tiny{[P1$\rangle$GM]}}
                \texttt{\{} \\
\texttt{  "current\_mode": "probing",} \\
\texttt{  "input\_rationale": "I have three examples: f({-}99, 10) = 723, f({-}1, {-}99) = {-}398, f({-}99, {-}2) = {-}501. The change in output with respect to b when a is constant ({-}99) suggests a coefficient of 102 for b. Specifically, 723 {-} ({-}501) = 1224, and 10 {-} ({-}2) = 12, so 1224/12 = 102. This suggests f(a, b) = 102*b + g(a). Let's test this hypothesis with the second example. If f({-}1, {-}99) = {-}398, then 102*({-}99) + g({-}1) = {-}398 => {-}10098 + g({-}1) = {-}398 => g({-}1) = 9700. Now let's check if this holds for a={-}99. f({-}99, 10) = 102*10 + g({-}99) = 723 => 1020 + g({-}99) = 723 => g({-}99) = {-}297} \\
            }
        }
    }
    & & \\ \\

    \theutterance \stepcounter{utterance}  
    & & & \multicolumn{4}{p{0.6\linewidth}}{
        \cellcolor[rgb]{0.9,0.9,0.9}{
            \makecell[{{p{\linewidth}}}]{
                \texttt{\tiny{[P1$\langle$GM]}}
                \texttt{Game Violation: ParseError: Response is not valid JSON.} \\
\texttt{Output ONLY:} \\
\texttt{One JSON with confidence\_score, current\_mode, input\_rationale, and hypothesis keys ...} \\
\texttt{OR} \\
\texttt{SOLVE: ```python ... ```} \\
            }
        }
    }
     \\ \\

    \theutterance \stepcounter{utterance}  
    & \multicolumn{4}{p{0.6\linewidth}}{
        \cellcolor[rgb]{0.561,0.737,0.561}{
            \makecell[{{p{\linewidth}}}]{
                \texttt{\tiny{[P1$\rangle$GM]}}
                \texttt{\{} \\
\texttt{  "current\_mode": "probing",} \\
\texttt{  "input\_rationale": "I have three examples: f({-}99, 10) = 723, f({-}1, {-}99) = {-}398, f({-}99, {-}2) = {-}501. The change in output with respect to b when a is constant ({-}99) suggests a coefficient of 102 for b. Specifically, 723 {-} ({-}501) = 1224, and 10 {-} ({-}2) = 12, so 1224/12 = 102. This suggests f(a, b) = 102*b + g(a). Let's test this hypothesis with the second example. If f({-}1, {-}99) = {-}398, then 102*({-}99) + g({-}1) = {-}398 => {-}10098 + g({-}1) = {-}398 => g({-}1) = 9700. Now let's check if this holds for a={-}99. f({-}99, 10) = 102*10 + g({-}99) = 723 => 1020 + g({-}99) = 723 => g({-}99) = {-}297} \\
            }
        }
    }
    & & \\ \\

    \theutterance \stepcounter{utterance}  
    & & & \multicolumn{4}{p{0.6\linewidth}}{
        \cellcolor[rgb]{0.9,0.9,0.9}{
            \makecell[{{p{\linewidth}}}]{
                \texttt{\tiny{[P1$\langle$GM]}}
                \texttt{Game Violation: ParseError: Response is not valid JSON.} \\
\texttt{Output ONLY:} \\
\texttt{One JSON with confidence\_score, current\_mode, input\_rationale, and hypothesis keys ...} \\
\texttt{OR} \\
\texttt{SOLVE: ```python ... ```} \\
            }
        }
    }
     \\ \\

    \theutterance \stepcounter{utterance}  
    & \multicolumn{4}{p{0.6\linewidth}}{
        \cellcolor[rgb]{0.561,0.737,0.561}{
            \makecell[{{p{\linewidth}}}]{
                \texttt{\tiny{[P1$\rangle$GM]}}
                \texttt{\{} \\
\texttt{  "current\_mode": "probing",} \\
\texttt{  "input\_rationale": "I have three examples: f({-}99, 10) = 723, f({-}1, {-}99) = {-}398, f({-}99, {-}2) = {-}501. The change in output with respect to b when a is constant ({-}99) suggests a coefficient of 102 for b. Specifically, 723 {-} ({-}501) = 1224, and 10 {-} ({-}2) = 12, so 1224/12 = 102. This suggests f(a, b) = 102*b + g(a). Let's test this hypothesis with the second example. If f({-}1, {-}99) = {-}398, then 102*({-}99) + g({-}1) = {-}398 => {-}10098 + g({-}1) = {-}398 => g({-}1) = 9700. Now let's check if this holds for a={-}99. f({-}99, 10) = 102*10 + g({-}99) = 723 => 1020 + g({-}99) = 723 => g({-}99) = {-}297} \\
            }
        }
    }
    & & \\ \\

    \theutterance \stepcounter{utterance}  
    & & & \multicolumn{4}{p{0.6\linewidth}}{
        \cellcolor[rgb]{0.9,0.9,0.9}{
            \makecell[{{p{\linewidth}}}]{
                \texttt{\tiny{[P1$\langle$GM]}}
                \texttt{Game Violation: ParseError: Response is not valid JSON.} \\
\texttt{Output ONLY:} \\
\texttt{One JSON with confidence\_score, current\_mode, input\_rationale, and hypothesis keys ...} \\
\texttt{OR} \\
\texttt{SOLVE: ```python ... ```} \\
            }
        }
    }
     \\ \\

    \theutterance \stepcounter{utterance}  
    & \multicolumn{4}{p{0.6\linewidth}}{
        \cellcolor[rgb]{0.561,0.737,0.561}{
            \makecell[{{p{\linewidth}}}]{
                \texttt{\tiny{[P1$\rangle$GM]}}
                \texttt{\{} \\
\texttt{  "current\_mode": "probing",} \\
\texttt{  "input\_rationale": "I have three examples: f({-}99, 10) = 723, f({-}1, {-}99) = {-}398, f({-}99, {-}2) = {-}501. The change in output with respect to b when a is constant ({-}99) suggests a coefficient of 102 for b. Specifically, 723 {-} ({-}501) = 1224, and 10 {-} ({-}2) = 12, so 1224/12 = 102. This suggests f(a, b) = 102*b + g(a). Let's test this hypothesis with the second example. If f({-}1, {-}99) = {-}398, then 102*({-}99) + g({-}1) = {-}398 => {-}10098 + g({-}1) = {-}398 => g({-}1) = 9700. Now let's check if this holds for a={-}99. f({-}99, 10) = 102*10 + g({-}99) = 723 => 1020 + g({-}99) = 723 => g({-}99) = {-}297} \\
            }
        }
    }
    & & \\ \\

    \theutterance \stepcounter{utterance}  
    & & & \multicolumn{4}{p{0.6\linewidth}}{
        \cellcolor[rgb]{0.9,0.9,0.9}{
            \makecell[{{p{\linewidth}}}]{
                \texttt{\tiny{[P1$\langle$GM]}}
                \texttt{Game Violation: ParseError: Response is not valid JSON.} \\
\texttt{Output ONLY:} \\
\texttt{One JSON with confidence\_score, current\_mode, input\_rationale, and hypothesis keys ...} \\
\texttt{OR} \\
\texttt{SOLVE: ```python ... ```} \\
            }
        }
    }
     \\ \\

    \theutterance \stepcounter{utterance}  
    & \multicolumn{4}{p{0.6\linewidth}}{
        \cellcolor[rgb]{0.561,0.737,0.561}{
            \makecell[{{p{\linewidth}}}]{
                \texttt{\tiny{[P1$\rangle$GM]}}
                \texttt{\{} \\
\texttt{  "current\_mode": "probing",} \\
\texttt{  "input\_rationale": "I have three examples: f({-}99, 10) = 723, f({-}1, {-}99) = {-}398, f({-}99, {-}2) = {-}501. The change in output with respect to b when a is constant ({-}99) suggests a coefficient of 102 for b. Specifically, 723 {-} ({-}501) = 1224, and 10 {-} ({-}2) = 12, so 1224/12 = 102. This suggests f(a, b) = 102*b + g(a). Let's test this hypothesis with the second example. If f({-}1, {-}99) = {-}398, then 102*({-}99) + g({-}1) = {-}398 => {-}10098 + g({-}1) = {-}398 => g({-}1) = 9700. Now let's check if this holds for a={-}99. f({-}99, 10) = 102*10 + g({-}99) = 723 => 1020 + g({-}99) = 723 => g({-}99) = {-}297} \\
            }
        }
    }
    & & \\ \\

    \theutterance \stepcounter{utterance}  
    & & & \multicolumn{4}{p{0.6\linewidth}}{
        \cellcolor[rgb]{0.9,0.9,0.9}{
            \makecell[{{p{\linewidth}}}]{
                \texttt{\tiny{[P1$\langle$GM]}}
                \texttt{Game Violation: ParseError: Response is not valid JSON.} \\
\texttt{Output ONLY:} \\
\texttt{One JSON with confidence\_score, current\_mode, input\_rationale, and hypothesis keys ...} \\
\texttt{OR} \\
\texttt{SOLVE: ```python ... ```} \\
            }
        }
    }
     \\ \\

    \theutterance \stepcounter{utterance}  
    & \multicolumn{4}{p{0.6\linewidth}}{
        \cellcolor[rgb]{0.561,0.737,0.561}{
            \makecell[{{p{\linewidth}}}]{
                \texttt{\tiny{[P1$\rangle$GM]}}
                \texttt{\{} \\
\texttt{  "current\_mode": "probing",} \\
\texttt{  "input\_rationale": "I have three examples: f({-}99, 10) = 723, f({-}1, {-}99) = {-}398, f({-}99, {-}2) = {-}501. The change in output with respect to b when a is constant ({-}99) suggests a coefficient of 102 for b. Specifically, 723 {-} ({-}501) = 1224, and 10 {-} ({-}2) = 12, so 1224/12 = 102. This suggests f(a, b) = 102*b + g(a). Let's test this hypothesis with the second example. If f({-}1, {-}99) = {-}398, then 102*({-}99) + g({-}1) = {-}398 => {-}10098 + g({-}1) = {-}398 => g({-}1) = 9700. Now let's check if this holds for a={-}99. f({-}99, 10) = 102*10 + g({-}99) = 723 => 1020 + g({-}99) = 723 => g({-}99) = {-}297} \\
            }
        }
    }
    & & \\ \\

    \theutterance \stepcounter{utterance}  
    & & & \multicolumn{4}{p{0.6\linewidth}}{
        \cellcolor[rgb]{0.9,0.9,0.9}{
            \makecell[{{p{\linewidth}}}]{
                \texttt{\tiny{[P1$\langle$GM]}}
                \texttt{Game Violation: ParseError: Response is not valid JSON.} \\
\texttt{Output ONLY:} \\
\texttt{One JSON with confidence\_score, current\_mode, input\_rationale, and hypothesis keys ...} \\
\texttt{OR} \\
\texttt{SOLVE: ```python ... ```} \\
            }
        }
    }
     \\ \\

    \theutterance \stepcounter{utterance}  
    & \multicolumn{4}{p{0.6\linewidth}}{
        \cellcolor[rgb]{0.561,0.737,0.561}{
            \makecell[{{p{\linewidth}}}]{
                \texttt{\tiny{[P1$\rangle$GM]}}
                \texttt{\{} \\
\texttt{  "current\_mode": "probing",} \\
\texttt{  "input\_rationale": "I have three examples: f({-}99, 10) = 723, f({-}1, {-}99) = {-}398, f({-}99, {-}2) = {-}501. The change in output with respect to b when a is constant ({-}99) suggests a coefficient of 102 for b. Specifically, 723 {-} ({-}501) = 1224, and 10 {-} ({-}2) = 12, so 1224/12 = 102. This suggests f(a, b) = 102*b + g(a). Let's test this hypothesis with the second example. If f({-}1, {-}99) = {-}398, then 102*({-}99) + g({-}1) = {-}398 => {-}10098 + g({-}1) = {-}398 => g({-}1) = 9700. Now let's check if this holds for a={-}99. f({-}99, 10) = 102*10 + g({-}99) = 723 => 1020 + g({-}99) = 723 => g({-}99) = {-}297} \\
            }
        }
    }
    & & \\ \\

    \theutterance \stepcounter{utterance}  
    & & & \multicolumn{4}{p{0.6\linewidth}}{
        \cellcolor[rgb]{0.9,0.9,0.9}{
            \makecell[{{p{\linewidth}}}]{
                \texttt{\tiny{[P1$\langle$GM]}}
                \texttt{Game Violation: ParseError: Response is not valid JSON.} \\
\texttt{Output ONLY:} \\
\texttt{One JSON with confidence\_score, current\_mode, input\_rationale, and hypothesis keys ...} \\
\texttt{OR} \\
\texttt{SOLVE: ```python ... ```} \\
            }
        }
    }
     \\ \\

    \theutterance \stepcounter{utterance}  
    & \multicolumn{4}{p{0.6\linewidth}}{
        \cellcolor[rgb]{0.561,0.737,0.561}{
            \makecell[{{p{\linewidth}}}]{
                \texttt{\tiny{[P1$\rangle$GM]}}
                \texttt{\{} \\
\texttt{  "current\_mode": "probing",} \\
\texttt{  "input\_rationale": "I have three examples: f({-}99, 10) = 723, f({-}1, {-}99) = {-}398, f({-}99, {-}2) = {-}501. The change in output with respect to b when a is constant ({-}99) suggests a coefficient of 102 for b. Specifically, 723 {-} ({-}501) = 1224, and 10 {-} ({-}2) = 12, so 1224/12 = 102. This suggests f(a, b) = 102*b + g(a). Let's test this hypothesis with the second example. If f({-}1, {-}99) = {-}398, then 102*({-}99) + g({-}1) = {-}398 => {-}10098 + g({-}1) = {-}398 => g({-}1) = 9700. Now let's check if this holds for a={-}99. f({-}99, 10) = 102*10 + g({-}99) = 723 => 1020 + g({-}99) = 723 => g({-}99) = {-}297} \\
            }
        }
    }
    & & \\ \\

    \theutterance \stepcounter{utterance}  
    & & & \multicolumn{2}{p{0.3\linewidth}}{
        \cellcolor[rgb]{0.9,0.9,0.9}{
            \makecell[{{p{\linewidth}}}]{
                \texttt{\tiny{[GM$|$GM]}}
                \texttt{Game Verdict: LOSS (turn limit reached)} \\
            }
        }
    }
    & & \\ \\

\end{supertabular}
}

\end{document}
